\documentclass[12pt,a4paper]{article}
\usepackage[utf8]{inputenc}
\usepackage[T2A]{fontenc}
\usepackage[russian]{babel}
\usepackage{amsmath,amssymb}
\usepackage{graphicx}
\usepackage{geometry}
\usepackage{setspace}
\usepackage{titlesec}
\usepackage{float}

\geometry{left=2cm,right=2cm,top=2cm,bottom=2cm}
\onehalfspacing

% \title{Лабораторная работа №4.6.2\\
% \large Туннелирование миллиметровых радиоволн}
% \author{}
% \date{}

\begin{document}

\title{\begin{center}Лабораторная работа № 4.6.2\end{center}
Туннелирование миллиметровых радиоволн}
\author{Струков О. И. \\ Б04-404}
\date{}

\thispagestyle{empty}
\pagenumbering{gobble}
\maketitle
\newpage
\pagenumbering{arabic}

\section*{Цель работы}
Экспериментальное исследование эффекта проникновения электромагнитных волн --- туннелирования --- через воздушный зазор между диэлектрическими призмами при полном внутреннем отражении на границе диэлектрик-воздух, а также моделирование интерферометра Майкельсона с использованием этого эффекта и измерение длины волны излучения и показателя преломления фторопласта для радиоволн миллиметрового диапазона.

\section*{Теоретическая часть}

\subsection*{Явление полного внутреннего отражения и туннелирования}
При падении электромагнитной волны на границу раздела двух сред с показателями преломления $n_1 > n_2$ под углом $\varphi_1$, превышающим предельный угол полного внутреннего отражения
\begin{equation}
\varphi_{\text{пред}} = \arcsin\left(\frac{n_2}{n_1}\right),
\end{equation}
волна полностью отражается от границы. Однако в среде с меньшим показателем преломления ($n_2$) возникает неоднородная (эванесцентная) волна, амплитуда которой экспоненциально убывает с расстоянием от границы:
\begin{equation}
E(z) \sim E_0 \exp\left(-\frac{z}{\Delta}\right),
\end{equation}
где $\Delta$ --- длина затухания, определяемая выражением
\begin{equation}
\Delta = \frac{\lambda_2}{2\pi\sqrt{n^2\sin^2\varphi_1 - 1}},
\label{eq:delta}
\end{equation}
$\lambda_2$ --- длина волны в воздухе, $n$ --- показатель преломления материала призмы, $\varphi_1$ --- угол падения.

Если на расстоянии, сравнимом с $\Delta$, от первой границы разместить вторую границу раздела (воздушный зазор между двумя призмами), то эванесцентная волна может возбудить распространяющуюся волну во второй призме. Это явление называется туннелированием электромагнитных волн.

\subsection*{Коэффициенты прохождения и отражения}
При туннелировании через воздушный зазор толщиной $\ell$ коэффициент прохождения $T$ убывает по экспоненциальному закону:
\begin{equation}
T \sim \exp\left(-\frac{2\ell}{\Delta}\right),
\label{eq:T}
\end{equation}
а коэффициент отражения $R$ связан с $T$ соотношением сохранения энергии:
\begin{equation}
T + R = 1.
\end{equation}

\subsection*{Интерферометр Майкельсона на основе туннелирования}
Воздушный зазор между призмами может использоваться в качестве делителя волны в схеме, моделирующей интерферометр Майкельсона. При перемещении подвижного зеркала на расстояние $\delta x$ изменяется разность хода интерферирующих лучей, что приводит к изменению интенсивности регистрируемого сигнала с периодом, равным половине длины волны излучения.

Для измерения показателя преломления материала пластинки интерференционным методом перед неподвижным зеркалом устанавливается пластина известной толщины $h$. При этом возникает приращение оптической длины пути:
\begin{equation}
\Delta L = 2h(n - 1).
\label{eq:optical_path}
\end{equation}
Это приращение компенсируется перемещением подвижного зеркала на расстояние $\delta x$, откуда показатель преломления определяется как:
\begin{equation}
n = 1 + \frac{\delta x}{h}.
\label{eq:n_interf}
\end{equation}
Для толстых пластин, когда $2h(n-1) > \lambda$, необходимо учитывать изменение порядка интерференции, используя приближённое значение $n \approx 1{,}5$.

\begin{figure}[H]
\centering
\includegraphics[width=0.7\textwidth]{1.jpg}
\caption{Схема установки для исследования туннелирования миллиметровых радиоволн}
\label{fig:scheme_tunneling}
\end{figure}

\begin{figure}[H]
\centering
\includegraphics[width=0.6\textwidth]{2.jpg}
\caption{Схема, моделирующая интерферометр Майкельсона}
\label{fig:scheme_michelson}
\end{figure}

\section*{Ход работы}

\subsection*{Подготовка приборов к работе}
\begin{enumerate}
    \item Генератор Г4--115 был настроен в соответствии с техническим описанием,
Столик с призмами был установлен так, чтобы воздушный зазор был ориентирован под углом $45^\circ$ к падающему лучу. Винт левого микрометра был установлен в положение, обеспечивающее максимальный диапазон изменения воздушного промежутка. Вращением винта правого микрометра воздушный промежуток был устранён.
Приёмная антенна была расположена на одной прямой с передатчиком. %Металлическое зеркало, стоящее на пути луча, было снято. Путём последовательных приближений, слегка поворачивая столик и приёмную антенну вокруг вертикальной оси, был достигнут максимальный отклик микроамперметра, после чего оба рейтера были закреплены.
    % \item Вращением ручек генератора (12 -- установка частоты и 7 -- установка генерации плавно) была настроена максимальная выходная мощность клистрона. При необходимости ток был уменьшен вращением ручки аттенюатора Ат на приёмной антенне. При отсутствии зазора прошедший сигнал составил $80\div90\%$ шкалы микроамперметра.
    % \item Вращением ручки 8 было достигнуто загорание контрольной лампочки 10. Была определена рабочая частота клистрона по шкале 9 и рассчитана соответствующая длина волны:
    \item Рабочая частота составила $\nu = 36$ ГГц, и отсюда была определена соответствующая длина волны:
    \[
    \lambda = \frac{c}{\nu} \approx 8,33 \text{ мм}.
    \]
% \end{enumerate}



\subsection*{Исследование туннелирования радиоволн}
% \begin{enumerate}
%     \setcounter{enumi}{5}
    \item Была снята зависимость интенсивности прошедшей волны от величины зазора $\ell$.% (использовался только правый микрометр, люфт устранялся). При увеличении зазора интенсивность падала монотонно, что свидетельствует о правильной ориентации призм.

После этого приёмник был переставлен для измерения отражённого сигнала и
% Путём небольшого поворота столика и приёмника было достигнуто равенство отклика микроамперметра на отражённый сигнал при максимальном зазоре и на прошедший сигнал при нулевом зазоре. 
была снята зависимость интенсивности отражённой волны от величины зазора.
    
    % \item Была установлена величина зазора, при которой ток равен половине максимального. Путём перестановки приёмника было подтверждено, что $T \approx R \approx 0{,}5$.
    
    \item На одном листе были построены графики зависимости коэффициентов $T$ и $R$ от величины зазора $\ell$, при этом токи были отнормированы на величину $I_{\text{max}}$. Заметно, что соотношение $T + R = 1$ приблизительно выполняется.
    
    \begin{figure}[H]
    \centering
    \includegraphics[width=0.9\textwidth]{1.pdf}
    \caption{Зависимости коэффициентов прохождения $T$ и отражения $R$ от толщины воздушного зазора $\ell$}
    \label{fig:graph_TR}
    \end{figure}
    
    \item Был построен график $\ln T = f(z)$, где $z$ -- показания микрометра. Полученные точки лежат на одной прямой согласно формуле~(\ref{eq:T}). По наклону прямой была рассчитана длина затухания $\Delta$: $k = (-0,455\pm0,024)$ 1/мм, $\Delta = \dfrac{-2}{k} = (4,4\pm0,2)$ мм, а затем по формуле~(\ref{eq:delta}) -- величина $n\sin\varphi_1 = (1,044\pm0,07)$. Поскольку $\varphi_1 \approx 45^{\circ}$, был найден показатель преломления фторопласта $n$: % (в условиях опыта не учитывалось, что входная плоскость призмы П1 наклонена на угол $8^\circ$ по отношению к фронту падающей волны):
    \[
    n = 1,48\pm0,10
    \]
    
    \begin{figure}[H]
    \centering
    \includegraphics[width=0.9\textwidth]{2.pdf}
    \caption{Зависимость $\ln T$ от показаний микрометра $z$}
    \label{fig:graph_lnT}
    \end{figure}
% \end{enumerate}

\subsection*{Интерферометр Майкельсона}
% \begin{enumerate}
%     \setcounter{enumi}{10}
    \item Была собрана схема интерферометра Майкельсона, в которой в качестве делителя использовался воздушный зазор между призмами. Оптимальный размер зазора соответствовал равенству $T \approx R \approx 0{,}5$. % На место было установлено неподвижное металлическое зеркало.
    
    \item Была снята зависимость тока от координаты $x$ подвижного зеркала. По графику $I = f(x)$ было определено экспериментальное значение длины волны электромагнитного излучения. %Поскольку максимумы были размыты, для повышения точности использовались координаты, соответствующие максимальным производным $dI/dx$. 
    Экспериментальная длина волны достаточно близка к величине, рассчитанной по рабочей частоте генератора:
    \[
    \lambda_{\text{эксп}} =  (8,6\pm0,2)\text{ мм}, \quad \lambda_{\text{расч}} = 8,33 \text{ мм}.
    \]
    
    \begin{figure}[H]
    \centering
    \includegraphics[width=0.9\textwidth]{3.pdf}
    \caption{Зависимость интенсивности сигнала от положения подвижного зеркала}
    \label{fig:graph_interference}
    \end{figure}
    
    \item Для измерения показателя преломления фторопласта интерференционным методом интерферометр был настроен на максимальную интенсивность, и перед неподвижным зеркалом была помещена пластина известной толщины $h = 6,2$ мм. Возникшее увеличение оптической длины пути было скомпенсировано отодвиганием подвижного зеркала от призм на расстояние $\delta x = 2,05$ мм. Показатель преломления фторопласта рассчитан по формуле~(\ref{eq:n_interf}):
    \[
    n = 1 + \frac{\delta x}{h} = (1,33\pm0,02)
    \]
    
    % \item Были сравнены результаты измерения показателя преломления фторопласта $n$ интерференционным методом и методом туннелирования.
\end{enumerate}



% \section*{Ответы на контрольные вопросы}

% \begin{enumerate}
%     \item \textbf{Какой угол называется предельным углом полного внутреннего отражения? Чему он равен?}
    
%     Предельным углом полного внутреннего отражения называется минимальный угол падения, при котором волна, падающая из оптически более плотной среды ($n_1$) на границу с оптически менее плотной средой ($n_2 < n_1$), полностью отражается. Он определяется выражением:
%     \[
%     \varphi_{\text{пред}} = \arcsin\left(\frac{n_2}{n_1}\right).
%     \]
%     Для границы фторопласт--воздух при $n \approx 1{,}5$ получаем $\varphi_{\text{пред}} \approx 41{,}8^\circ$.
    
%     \item \textbf{Что такое неоднородная волна? Как направлен вектор Умова--Пойнтинга в такой волне и почему?}
    
%     Неоднородной (эванесцентной) называется волна, амплитуда которой экспоненциально убывает в направлении, перпендикулярном границе раздела сред, при этом фазовые фронты параллельны границе. В такой волне вектор Умова--Пойнтинга направлен вдоль границы раздела сред, поскольку перенос энергии в направлении нормали к границе отсутствует --- энергия не проникает во вторую среду, а лишь ``скользит'' вдоль границы.
    
%     \item \textbf{В чём заключается явление туннелирования волн? Чем отличается это явление от обычного прохождения волн через тонкие прозрачные пластинки?}
    
%     Явление туннелирования заключается в проникновении электромагнитной волны через слой среды, в котором распространение волны в обычном смысле невозможно (например, через воздушный зазор при углах падения больше предельного). В отличие от обычного прохождения через прозрачную пластинку, при туннелировании:
%     \begin{itemize}
%         \item в промежуточной среде существует не распространяющаяся, а затухающая эванесцентная волна;
%         \item коэффициент прохождения убывает экспоненциально с увеличением толщины зазора;
%         \item явление наблюдается только при выполнении условий полного внутреннего отражения на первой границе.
%     \end{itemize}
    
%     \item \textbf{Каковы возможные причины расхождения измеренных в работе величин друг с другом и с табличными значениями этих величин?}
    
%     Возможные причины расхождений:
%     \begin{itemize}
%         \item неточность установки угла падения $45^\circ$ и юстировки оптической схемы;
%         \item наличие люфтов в микрометрических винтах;
%         \item конечная точность отсчёта показаний микроамперметра;
%         \item неидеальность поверхностей призм и наличие загрязнений;
%         \item влияние обратных отражений от элементов установки;
%         \item приближённость теоретических формул (неучёт наклона входной грани призмы на $8^\circ$, дисперсии материала);
%         \item погрешность определения рабочей частоты генератора.
%     \end{itemize}
% \end{enumerate}

\section*{Вывод}

В ходе выполнения лабораторной работы было экспериментально исследовано явление туннелирования миллиметровых радиоволн через воздушный зазор между фторопластовыми призмами при полном внутреннем отражении. Были сняты зависимости коэффициентов прохождения и отражения от толщины зазора, подтверждено выполнение закона сохранения энергии $T + R = 1$.% и экспоненциальный характер убывания прошедшего сигнала. 
По наклону зависимости $\ln T(\ell)$ была определена длина затухания волны и рассчитан показатель преломления фторопласта: $n_{\text{тун}} = (1,48\pm0,10)$.

Была собрана и исследована схема интерферометра Майкельсона,% в которой воздушный зазор между призмами использовался в качестве делителя волны. 
по периоду интерференционной зависимости интенсивности от положения подвижного зеркала была определена длина волны излучения: $\lambda_{\text{эксп}} = (8,6\pm0,2)$ мм, что близко к расчётному значению по частоте генератора. Интерференционным методом с использованием пластины известной толщины был измерен показатель преломления фторопласта: $n_{\text{инт}} = (1,33\pm0,02)$.

Результаты измерений показателя преломления двумя независимыми методами оказались близки друг к другу и к табличному значению для фторопласта ($n \approx 1{,}4\div1{,}5$). Наиболее точным оказался метод расчёта через коэффициент туннелирования.%, что подтверждает корректность проведённых измерений и применимость теоретических моделей туннелирования и интерференции электромагнитных волн миллиметрового диапазона.

\end{document}